\chapter{Generalidades}

\section{Planteamiento del problema}

\subsection{Descripción del problema}

La región Cusco enfrenta un incremento sostenido en la percepción de inseguridad ciudadana y en la ocurrencia de hechos delictivos. Según registros de la \gls{pnp}, las denuncias por homicidio en el departamento de Cusco aumentaron de 36 casos en 2022 a 60 casos en 2024 (enero-noviembre) \citep{pnp2024homicidios}. Asimismo, la \gls{enapres} ubicó a Cusco entre los departamentos con mayor porcentaje de percepción de inseguridad del país, con un 92,3\,\% durante el periodo noviembre 2022-abril 2023 \citep{inei2023enapres}.

La literatura criminológica documenta que el delito no se reparte al azar: una proporción reducida de lugares concentra una parte desproporcionada de los hechos \citep{sherman1989,weisburd2015}. Por ello, los gobiernos locales, que participan en el Sistema Nacional de Seguridad Ciudadana a través de los comités distritales \citep{ley27933}, necesitan saber en qué zonas y en qué momentos se concentran los incidentes para orientar el patrullaje.

El Mapa del Delito Georreferenciado del Ministerio del Interior publica una capa denominada SIPCOP Municipal con registros de ocurrencias del patrullaje municipal e integrado (\termref{term-serenazgo}{serenazgo} y patrullaje local integrado). El extracto disponible para el distrito de San Sebastián contiene 996 registros entre el 1 de enero y el 29 de agosto de 2025, con categoría, turno y coordenadas. Estos datos ofrecen un caso para especificar y evaluar una \termref{term-architecture}{arquitectura analítica} que integre su validación, procesamiento espacial, agrupamiento y presentación de resultados. El campo temporal disponible es la fecha de alerta; no se cuenta con hora individual en el archivo descargado. Las coordenadas están en grados decimales con tres decimales (unos 110 m), de modo que los 996 registros se reducen a 412 ubicaciones distintas.

\subsection{Identificación del problema}

La principal dificultad es la ausencia de una arquitectura documentada y reproducible que convierta los registros disponibles de SIPCOP Municipal en resultados espaciales y temporales trazables para la gestión local. Un análisis aislado de los datos no resuelve cómo integrar la recepción y validación de registros, la preparación de coordenadas, la selección y comparación de algoritmos, la evaluación estadística de las concentraciones y la generación de salidas cartográficas. La arquitectura debe además considerar las limitaciones del extracto: 996 registros fechados entre enero y agosto de 2025, sin hora individual y con coordenadas de resolución aproximada de 110 m. El periodo no representa los años 2022-2024 ni los últimos meses de 2025.

\subsection{Formulación del problema}

\subsubsection{Problema general}

¿Cómo debe estructurarse y evaluarse una arquitectura basada en clústeres espaciales para caracterizar los incidentes de seguridad ciudadana registrados en SIPCOP Municipal en San Sebastián, Cusco?

\subsubsection{Problemas específicos}

\begin{enumerate}[label=PE\arabic*.]
    \item ¿Qué requisitos funcionales, de datos y de calidad debe atender la arquitectura, dadas las características del extracto SIPCOP-M?
    \item ¿Qué componentes, interfaces y flujo de procesamiento deben conformar la arquitectura para integrar la validación de datos, el agrupamiento espacial y el análisis temporal?
    \item ¿Cómo implementar una prueba de concepto que ejecute el flujo arquitectónico con técnicas de agrupamiento espacial y produzca resultados reproducibles?
    \item ¿En qué medida la arquitectura y su prueba de concepto cumplen los requisitos funcionales y los criterios de calidad definidos?
\end{enumerate}

\section{Objetivos de la investigación}
\label{sec:objetivos}

\subsection{Objetivo general}

Desarrollar y evaluar una arquitectura analítica basada en clústeres espaciales para caracterizar los incidentes de seguridad ciudadana registrados en SIPCOP Municipal en San Sebastián, Cusco, mediante una \termref{term-proof-of-concept}{prueba de concepto} reproducible aplicada al extracto disponible.

\subsection{Objetivos específicos}

\begin{enumerate}[label=OE\arabic*.]
    \item Analizar las características y calidad del extracto SIPCOP-M y especificar los requisitos funcionales, de datos y de calidad de la arquitectura.
    \item Diseñar la arquitectura, definiendo sus componentes, interfaces, flujo de datos y mecanismos de trazabilidad para el procesamiento espacial y temporal.
    \item Implementar una prueba de concepto que integre la preparación de datos, DBSCAN, la comparación con HDBSCAN, los contrastes estadísticos y la generación de resultados cartográficos.
    \item Evaluar la prueba de concepto frente a los requisitos y criterios de aceptación definidos, considerando funcionalidad, reproducibilidad, calidad de resultados y limitaciones del conjunto de datos.
\end{enumerate}

\section{Hipótesis de la investigación}
\label{sec:hipotesis}

\subsection{Hipótesis general}

\textbf{H\textsubscript{G}:} La arquitectura analítica propuesta, implementada como prueba de concepto, permitirá ejecutar de forma trazable y reproducible el flujo de validación, agrupamiento espacial, análisis temporal y presentación de resultados para el extracto SIPCOP-M de San Sebastián, de acuerdo con los requisitos y criterios de aceptación establecidos.

\subsection{Hipótesis específicas}

\begin{enumerate}[label=H\textsubscript{\arabic*}:]
    \item[H\textsubscript{1}:] Los requisitos funcionales, de datos y de calidad definidos podrán vincularse con componentes y pruebas de aceptación explícitos en una matriz de trazabilidad completa.
    \item[H\textsubscript{2}:] La arquitectura especificará todos los componentes, interfaces y transformaciones necesarios para ejecutar el flujo desde la entrada del extracto hasta la generación de resultados espaciales y temporales.
    \item[H\textsubscript{3}:] La prueba de concepto procesará el extracto conforme a las reglas de validación, transformación y trazabilidad especificadas, y generará resultados reproducibles al repetir la ejecución con los mismos datos, parámetros y configuración.
    \item[H\textsubscript{4}:] La evaluación de la prueba de concepto demostrará el cumplimiento de los criterios de aceptación definidos antes de la ejecución, sin exclusiones de registros ni parámetros no documentados.
\end{enumerate}

\noindent Las hipótesis se evaluarán mediante pruebas funcionales, repetición controlada del flujo y verificación de los productos generados; la existencia de concentraciones significativas no se presupone y constituye un resultado del caso de aplicación, no una condición de éxito de la arquitectura.

\section{Justificación de la investigación}
\label{sec:justificacion}

\textbf{Justificación científica.} La investigación aportará una arquitectura analítica reproducible para integrar métodos de agrupamiento espacial, evaluación estadística y caracterización temporal en registros administrativos de seguridad ciudadana, documentando sus supuestos y límites mediante un caso local.

\textbf{Justificación metodológica.} El estudio especificará componentes, interfaces, flujo de datos, reglas de trazabilidad y criterios de aceptación, y probará el flujo con un artefacto reproducible. Así se podrá distinguir la calidad de la arquitectura de los resultados particulares obtenidos en el extracto.

\textbf{Justificación tecnológica.} La prueba de concepto integrará herramientas de código abierto para convertir datos tabulares del SIPCOP-M en resultados espaciales y temporales verificables. La arquitectura será independiente de una plataforma institucional específica y podrá servir de referencia para futuras implementaciones, sin afirmar que constituye un sistema de producción.

\textbf{Justificación social y práctica.} La Municipalidad Distrital dispondrá de información sobre zonas críticas y momentos de mayor recurrencia para focalizar el patrullaje y la prevención del serenazgo. El estudio no mide el efecto de ninguna intervención sobre la reducción del delito.

\section{Alcances y limitaciones}

\subsection{Alcances}

\textbf{Alcance espacial.} Distrito de San Sebastián, provincia y departamento de Cusco.

\textbf{Alcance temporal.} Registros con fecha de alerta entre el 1 de enero y el 29 de agosto de 2025, según el extracto disponible al preparar el estudio.

\textbf{Alcance tecnológico.} Especificación de una arquitectura analítica lógica; prototipo de prueba de concepto con ingesta local de archivos, validación y preparación de datos, agrupamiento con DBSCAN y comparación con HDBSCAN, contraste espacial y análisis temporal, generación de salidas tabulares y cartográficas, trazabilidad de parámetros y pruebas de aceptación. No incluye despliegue institucional ni integración en producción.

\textbf{Alcance disciplinar.} Ciencia de datos y análisis geoespacial, dentro de la Ingeniería Informática y de Sistemas.

\subsection{Limitaciones}

\textbf{Calidad y cobertura del registro.} El análisis se limita a 996 filas del extracto público descargado. Se evaluarán coordenadas, categorías, duplicados y otros problemas de calidad; los resultados no cubren registros no incluidos en ese extracto. El archivo no trae un identificador único de ocurrencia, por lo que los duplicados no pueden confirmarse, y todos los registros pertenecen a un mismo genérico (alertas tempranas en apoyo a la PNP), por lo que el extracto parece ser un subconjunto de la actividad total.

\textbf{Resolución espacial.} Las coordenadas tienen tres decimales (unos 110 m): hay 412 ubicaciones distintas y el 77\,\% de los registros comparte coordenada exacta con otro. Es probable que parte de las ubicaciones proceda de una geocodificación por puntos de referencia (por ejemplo, paraderos o el cementerio) y no de la posición real del hecho; los resultados no se interpretarán a escalas menores que esa resolución.

\textbf{Estabilidad del registro en el tiempo.} Los registros por día varían de 2,0 (junio) a 7,5 (agosto); hay 26 de 241 días sin ningún registro (no se sabe si hubo cero eventos o falta de reporte) y la modalidad de ``exposición a peligro o abandono de personas'' pasa del 57\,\% de los registros en enero a 9--17\,\% de junio a agosto, lo que sugiere un cambio de clasificación o de práctica de registro. Las diferencias entre meses pueden reflejar el registro y no el fenómeno.

\textbf{Cobertura temporal.} El conjunto disponible abarca del 1 de enero al 29 de agosto de 2025. No permite describir 2022-2024, el último cuatrimestre de 2025 ni inferir un patrón anual completo.

\textbf{Resolución temporal.} El archivo contiene fecha de alerta y turno, pero no hora individual de ocurrencia; por ello no se analizarán patrones horarios y la fecha de alerta puede no coincidir exactamente con la del hecho.

\textbf{Registro de lo atendido, no de lo ocurrido.} La capa SIPCOP Municipal documenta ocurrencias registradas por el patrullaje municipal e integrado (el 51\,\% de los registros proviene de Patrullaje Local Integrado), que incluyen hechos de seguridad que no son necesariamente delitos (el 34\,\% corresponde a exposición a peligro o abandono de personas); no incluye los hechos no reportados ni los denunciados solo ante la Policía Nacional. Los patrones describen lo registrado en la capa y no la totalidad del delito ni necesariamente todos los servicios atendidos.

\textbf{Precisión de la ubicación.} La exactitud de las coordenadas depende del modo de registro y condiciona la elección del radio $\varepsilon$.

\textbf{Sensibilidad a los parámetros.} Los grupos obtenidos con DBSCAN dependen de $\varepsilon$ y MinPts; se documentará un análisis de sensibilidad.

\textbf{Alcance inferencial.} El diseño es observacional: los resultados describen patrones y no establecen causas ni son generalizables a otros distritos. El registro refleja también dónde y cuándo se patrulla: el 76\,\% de los registros nace del patrullaje y solo 232 proceden de llamadas ciudadanas, alarmas o mensajería, lo que puede inflar la presencia de ocurrencias en zonas y turnos con mayor cobertura. Este sesgo se mide con un análisis de sensibilidad por origen del registro.
